% Source: https://artofproblemsolving.com/community/c3h3630844_partial_theory_1perpendicular_hypotenuse
% Stable identity: aops-topic-3630844-q01
Pythagoras gives $b^2+p^2=(p+1)^2$, so $b^2=2p+1$. Thus $b$ is odd; write $b=2k+1$. Then $p=2k(k+1)$ and the hypotenuse is $2k(k+1)+1$, for $k\ge 1$. Conversely these lengths satisfy Pythagoras and $b\le p$, so they are all solutions. The legs $1,2$ instead give hypotenuse $\sqrt 5$, not $3$, because $1^2+2^2=5\ne 9$.
\begin{asy}
import geometry; size(220);
pair A=(0,0),B=(4,0),C=(0,3);draw(A--B--C--cycle);draw((0.2,0)--(0.2,0.2)--(0,0.2));label("$A$",A,SW);label("$B$",B,SE);label("$C$",C,NW);label("$p$",(2,0),S);label("$b$",(0,1.5),W);label("$p+1$",(2,1.5),NE);
\end{asy}
